% Generated draft: requires mathematical and source review.
\documentclass[11pt]{article}
\usepackage[margin=25mm]{geometry}
\usepackage{fontspec}
\setmainfont{Latin Modern Roman}
\usepackage{amsmath,amssymb,mathtools,unicode-math}
\setmathfont{Latin Modern Math}
\usepackage{enumitem,needspace,xurl,hyperref}
\setlength{\parindent}{0pt}
\setlength{\parskip}{5pt}
\setlength{\emergencystretch}{3em}
\newcommand{\N}{\mathbb{N}}
\newcommand{\Z}{\mathbb{Z}}
\newcommand{\Q}{\mathbb{Q}}
\newcommand{\R}{\mathbb{R}}
\newcommand{\C}{\mathbb{C}}
\newcommand{\F}{\mathbb{F}}
\newcommand{\A}{\mathbb{A}}
\newcommand{\PP}{\mathbb{P}}
\newcommand{\E}{\mathbb{E}}
\newcommand{\eps}{\varepsilon}
\newcommand{\dd}{\,\mathrm{d}}
\begin{document}
\section*{UnsolvedMath ID: 30001006}
\textbf{Source problem code:} catalogue entry
\subsection{Definitions and mathematical statement}
Let $\mathcal A_n$ be the space of algebraic curvature operators on $\R^n$, namely self-adjoint endomorphisms of $\Lambda^2\R^n$ satisfying the first Bianchi identity. Write $\operatorname{scal}(R)$ for scalar curvature and $R_W$ for the Weyl component. For the source's norm and scalar-curvature normalization define \[\mathcal C_d=\{R\in\mathcal A_n:\operatorname{scal}(R)>d\|R_W\|\},\qquad d(n)=\sqrt{2(n-2)(n-1)}.\] The parameter $d$ in the family is distinct from the dimension $n$; $d(n)$ is the distinguished dimension-dependent value supplied in the statement. The norm is not identified as an operator norm or a tensor/Hilbert--Schmidt norm, so the numerical threshold cannot be interpreted by silently choosing one. A geometric meaning of Ricci-flow invariance is that a smooth Ricci flow $\partial_tg=-2\operatorname{Ric}(g)$ with initial curvature operators in $\mathcal C_d$ retains that pointwise condition while the flow exists, under appropriate global hypotheses. The extract does not specify these hypotheses or whether its theorem instead formulates invariance for the algebraic curvature reaction ODE, nor whether a closed cone or its strict interior is the actual preservation object.\par The source asks which values of $d$ yield invariance. Its original wording asks more particularly whether invariance is possible only at the values singled out by the report's theorem. Those singled-out values and the theorem's hypotheses are absent, so no asserted list is supplied here. The second question asks whether the report's high-dimensional invariance theorem can be sharpened to threshold $n_0=12$, that is, apply in every dimension $n\geq12$ with the same intended curvature-condition hypotheses. This cannot be replaced by an invented assertion about all $d$ or solely about $d(n)$ without the missing theorem. The supplied 2026-08-21 triage labels the record partially solved, reports high-dimensional invariant-cone results in the report, and says neither the full classification nor the $n_0=12$ claim is settled; it describes later work as further invariant conditions and quantitative evidence rather than a full proof.
\subsection{Short English statement}
Classify the Ricci-flow-preserved scalar-versus-Weyl curvature cones and ask whether the report's high-dimensional theorem already holds from dimension twelve. The norm normalization, distinguished parameter list, and theorem hypotheses are missing.
\subsection{Sources}
\begin{itemize}\item Local UnsolvedMath registry, permanent ID 30001006, source problem catalogue entry, \emph{Sharp Ricci-Flow-Invariant Curvature Cones}, proposed year 2008; no individual proposer supplied.\item Supplied citation: \emph{Geometrie} (2009), Oberwolfach Reports, DOI 10.4171/owr/2008/34, pp.\ 1933--1978: \url{https://doi.org/10.4171/owr/2008/34}. Supplied literature label: original Oberwolfach report.\item Supplied literature label: recent curvature-condition evidence, \url{https://arxiv.org/abs/2412.13633}.\item Supplied literature label: noncoercive invariant-cone results, \url{https://arxiv.org/abs/1308.1190}.\item Literature assessment date: 2026-08-21. The status and progress claims are supplied dated curation, not independently verified here.\end{itemize}
\end{document}
